\documentclass[11pt,a4paper]{article}

\usepackage[margin=2.6cm]{geometry}
\usepackage{amsmath,amssymb,amsthm,mathtools}
\usepackage{microtype}
\usepackage{booktabs,array,tabularx}
\usepackage[shortlabels]{enumitem}
\usepackage{xcolor}
\usepackage[colorlinks=true,linkcolor=black,citecolor=black,urlcolor=blue!50!black]{hyperref}

\setlist{itemsep=2pt,topsep=4pt}
\setlength{\parskip}{3pt}

\theoremstyle{plain}
\newtheorem{theorem}{Theorem}
\newtheorem{lemma}[theorem]{Lemma}
\newtheorem{proposition}[theorem]{Proposition}
\theoremstyle{definition}
\newtheorem{definition}[theorem]{Definition}
\theoremstyle{remark}
\newtheorem{remark}[theorem]{Remark}

\newcommand{\F}{\mathbb{F}_2}
\newcommand{\Aff}{\mathrm{Aff}}
\newcommand{\Span}{\operatorname{span}}
\newcommand{\supp}{\operatorname{supp}}
\newcommand{\Lam}{\Lambda^{2}}
\newcommand{\Dir}{E}
\newcommand{\lean}[1]{\texttt{\detokenize{#1}}}

\title{\bfseries An alternative proof of Theorem 12}
\author{AI-written draft, pending human review}
\date{2 November 2026}

\begin{document}
\maketitle

\begin{abstract}
\noindent
We give a second proof of Theorem~12 of \texttt{lambda4\_proof.tex} (cited as [L4]): no circuit with five AND gates, four of them linear, computes the three low bits $(z_0,z_1,z_2)$ of the product of two 3-bit integers. The proof replaces Lemma~5, Proposition~7, the dimension count of Lemma~4 and the reduction to $\dim Q=4$ by one strengthening of the derivative lemma, in which the contraction $\iota_e(Q)$ is allowed to be one-dimensional. The case split is on how many of the three linear gates occur in a decomposition of $Z_1$, and uses only Lemmas~3, 6 and~11 of [L4]. The statement proved is slightly stronger than Theorem~12: the three forms $P_1,P_2,P_3$ need only have rank at most $2$, and $\dim Q$ is arbitrary. Every step below is a theorem of the Lean development \texttt{lowhalf-mul3}, checked by the Lean kernel with no axioms beyond \lean{propext}, \lean{Classical.choice} and \lean{Quot.sound}; Section~\ref{sec:table} lists the correspondence.
\end{abstract}

\section{Setting}

We use the notation of [L4] throughout. Functions on $\F^6$ are identified with their algebraic normal forms, $\Aff=\Span\{1,x_1,\dots,x_6\}$, and for $\deg f\le2$ we write $f=f_0+f_1+f_2$ and identify $f_2$ with a 2-form. The inputs are $a_0,a_1,a_2,b_0,b_1,b_2$, and
\[
  z_0=a_0b_0,\quad z_1=a_0b_1+a_1b_0,\quad z_2=q_2+a_0a_1b_0b_1,\quad q_2=a_0b_2+a_1b_1+a_2b_0,
\]
\[
  P_0=a_0\wedge b_0,\qquad Z_1=a_0\wedge b_1+a_1\wedge b_0,\qquad U=\langle a_0,a_1,b_0,b_1\rangle .
\]
For a space $Q$ of 2-forms, $V_4(Q)=\{f:\deg f\le2,\ f_2\in Q\}$, and the condition on the last gate is
\begin{equation}\label{eq:star}
  f,h\in V_4(Q),\qquad fh\in z_2+V_4(Q).
\end{equation}
The directions are $\Dir=\{e_{a_2},e_{b_2},e_{a_2}+e_{b_2}\}$, $\iota_e$ is contraction with $e$, and $\delta_e=\iota_e(q_2)$, which is $b_0$, $a_0$ or $a_0+b_0$. We record the cross part of a 2-form $\omega$ as the pair of contractions $\Phi(\omega)=(\iota_{e_{a_2}}\omega,\ \iota_{e_{b_2}}\omega)$; this carries the same information as the triple $(\alpha,\beta,\gamma)$ of [L4], with $\gamma$ the $b_2$-coefficient of $\iota_{e_{a_2}}\omega$. In particular $\Phi(P_0)=\Phi(Z_1)=0$ and $\iota_e(P_0)=0$ for every $e\in\Dir$.

For a linear form $x$ let $\pi(x)=(x(e_{a_2}),x(e_{b_2}))\in\F^2$. By Lemma~11 of [L4] (Types), every form $l\wedge r$ can be rewritten as $l'\wedge r'$ with $\langle l',r'\rangle=\langle l,r\rangle$ in one of three normal forms: \emph{cross part zero} ($\pi(l')=\pi(r')=0$); \emph{type $v$} ($\pi(r')=0$, i.e.\ $r'\in U$, and $\pi(l')=v\ne0$); or \emph{full} ($l'=a_2+x$, $r'=b_2+y$ with $x,y\in U$). For type $v$ we have $\iota_e(l'\wedge r')=l'(e)\,r'$; for the full type $\iota_{e_{a_2}}(l'\wedge r')=r'$ and $\iota_{e_{b_2}}(l'\wedge r')=l'$. The $a_2\wedge b_2$-coefficient $\gamma$ is $0$ for type $v$ and $1$ for the full type.

\paragraph{Two facts about $Z=Z_1+c_0P_0$.} For $c_0\in\F$ the form $Z=Z_1+c_0P_0$ lies in $\Lam(U)$, so $\Phi(Z)=0$, and it has rank $4$: its Pfaffian on $a_0,a_1,b_0,b_1$ is $1$. Hence $Z$ is not of the form $l\wedge r$, since every $l\wedge r$ has all its $4\times4$ Pfaffians equal to $0$.

\section{The derivative lemma with a one-dimensional contraction}

For a function $p$ of degree at most $2$ and $e\in\Dir$, $D_ep(x)=p(x+e)+p(x)$ is affine with linear part $\iota_e(p_2)$:
\begin{equation}\label{eq:bridge}
  D_ep(x)=D_ep(0)+\langle\iota_e(p_2),x\rangle .
\end{equation}
This is the first line of the proof of [L4, Lemma~8].

\begin{lemma}[Lemma~8$'$]\label{lem:8p}
Let $e\in\Dir$, and suppose $\iota_e(Q)\subseteq\langle\ell\rangle$ for one linear form $\ell\ne\delta_e$. Then \eqref{eq:star} has no solution.
\end{lemma}

The case $\ell=0$ is [L4, Lemma~8].

\begin{proof}
By \eqref{eq:bridge}, every $p\in V_4(Q)$ has $D_ep=c_p+\gamma_p\ell$ with $c_p,\gamma_p\in\F$. Let $(f,h)$ be a solution and $g=fh+z_2\in V_4(Q)$.

\emph{Normalisation.} We may assume $\gamma_f=0$. If $\gamma_f=1$ and $\gamma_h=0$, exchange $f$ and $h$. If $\gamma_f=\gamma_h=1$, replace $(f,h)$ by $(f+h,h)$: since $h^2=h$, $(f+h)h+z_2=g+h\in V_4(Q)$, and $\gamma_{f+h}=0$.

So $D_ef=c\in\F$. Put $H=D_eh=c_h+\gamma_h\ell$ and $K=D_eg=c_g+\gamma_g\ell$, both affine. As in [L4, Lemma~8], $D_ez_2=\delta_e$ and the product rule $D_e(fh)=D_ef\cdot h+f\cdot D_eh+D_ef\cdot D_eh$ give, pointwise,
\[
  K=c\,h+H\,f+c\,H+\delta_e .
\]
\emph{Case $c=1$.} Then $h=H(f+1)+A$ with $A=K+\delta_e$ affine, so $fh=H\,f(f+1)+fA=fA$ has degree at most $3$.

\emph{Case $c=0$.} Then $Hf=A$ with $A=K+\delta_e=c_g+\gamma_g\ell+\delta_e$. We claim $H\equiv1$. Otherwise $H$ vanishes somewhere, and $A$ vanishes on the zero set of $H$. If $H$ is constant, then $H\equiv0$ and $A\equiv0$, so $\delta_e=c_g+\gamma_g\ell$. Otherwise $H=c_h+\ell$ with $\ell\ne0$, and $A$ is an affine function vanishing on the hyperplane $\{\ell=c_h\}$, so $A\in\{0,\ell+c_h\}$. In both cases $\delta_e\in\Span\{1,\ell\}$, which is false because $\delta_e$ is a nonzero linear form different from $\ell$. Hence $H\equiv1$, $f=A$ is affine, and $\deg(fh)\le3$.

In both cases $fh=g+z_2$ with $\deg g\le2$, so the degree-4 part of $fh$ is $a_0a_1b_0b_1\ne0$, a contradiction.
\end{proof}

\begin{lemma}[Case A]\label{lem:A}
Let $Q=\Span\{P_0,\dots,P_3\}$, and suppose there is a 2-form $R=a\wedge b$ such that $\Phi(P_k)\in\{0,\Phi(R)\}$ for every $k$. Then \eqref{eq:star} has no solution.
\end{lemma}

\begin{proof}
By [L4, Lemma~6], $\Phi(R)\ne\Phi(q_2)=(b_0,a_0)$ (for $R=0$ this is clear, as $\Phi(q_2)\ne0$). So $\iota_{e_{a_2}}R\ne b_0=\delta_{e_{a_2}}$ or $\iota_{e_{b_2}}R\ne a_0=\delta_{e_{b_2}}$. In the first case put $e=e_{a_2}$, otherwise $e=e_{b_2}$, and $\ell=\iota_eR\ne\delta_e$. Each $\iota_e(P_k)$ is $0$ or $\ell$, so $\iota_e(Q)\subseteq\langle\ell\rangle$, and Lemma~\ref{lem:8p} applies.
\end{proof}

\section{The theorem}

\begin{theorem}[Theorem~12 of {[L4]}, strengthened]\label{thm:main}
Let $P_1,P_2,P_3$ be 2-forms of rank at most $2$, $Q=\Span\{P_0,P_1,P_2,P_3\}$, and $Z_1\in Q$. Then \eqref{eq:star} has no solution. Consequently, no circuit with five AND gates, four of them linear, computes $(z_0,z_1,z_2)$.
\end{theorem}

\begin{proof}
Write $P_k=l_k\wedge r_k$, and $Z_1=c_0P_0+c_1P_1+c_2P_2+c_3P_3$. Put $Z=Z_1+c_0P_0$ and $S=\{k\ge1:c_k=1\}$, so that $\sum_{k\in S}P_k=Z$. Taking cross parts gives $\sum_{k\in S}\Phi(P_k)=\Phi(Z)=0$.

\medskip\noindent\emph{$|S|\le1$.} Then $Z=0$ or $Z=P_k$ is of the form $l\wedge r$, which is impossible.

\medskip\noindent\emph{$|S|=2$, say $S=\{j,k\}$.} Then $P_j+P_k=Z$. Neither summand is $0$, since $Z$ itself is not of the form $l\wedge r$, so both have rank $2$, and Lemma~3 of [L4] puts both in $\Lam(U)$: $\Phi(P_j)=\Phi(P_k)=0$. For the remaining index $m$, Case~A applies with $R=P_m$.

\medskip\noindent\emph{$|S|=3$.} Then $P_1+P_2+P_3=Z$ and $\Phi(P_1)+\Phi(P_2)+\Phi(P_3)=0$. If some $\Phi(P_k)=0$, the other two cross parts are equal, and Case~A applies with $R$ either of them. So assume all three cross parts are nonzero, and bring each $P_k$ to its normal form (Lemma~11); none has cross part zero.

\emph{Parity.} The $a_2\wedge b_2$-coefficients satisfy $\gamma(P_1)+\gamma(P_2)+\gamma(P_3)=\gamma(Z)=0$, and each is $0$ for type $v$ and $1$ for the full type. So either none or exactly two of the $P_k$ are full.

\emph{None full.} Then $P_k=l_k\wedge r_k$ with $r_k\in U$ and $v_k=\pi(l_k)\ne0$, and $\iota_e(P_k)=v_k(e)\,r_k$ for $e\in\Dir$. The cross-part relation gives
\[
  \textstyle\sum_k v_k(e)\,r_k=0\qquad\text{for every }e\in\Dir .
\]
Here $v\mapsto v(e)$ runs over the three nonzero linear functionals on $\F^2$, and each of them vanishes on exactly one nonzero vector.
\begin{enumerate}[(a)]
  \item If $v_1=v_2=v_3=v$, take $e$ with $v(e)=0$. Then $\iota_e(P_k)=0$ for all $k$, and also $\iota_e(P_0)=0$, so $\iota_e(Q)=0$ and [L4, Lemma~8] applies.
  \item If exactly two agree, say $v_i=v_j\ne v_m$, take $e$ with $v_i(e)=0$. Then $v_m(e)=1$, and the relation gives $r_m=0$, so $\Phi(P_m)=0$, contrary to the assumption.
  \item If $v_1,v_2,v_3$ are distinct, they are the three nonzero vectors of $\F^2$. For each $i$, the $e$ with $v_i(e)=0$ has $v_j(e)=v_k(e)=1$ for the other two indices, so $r_j=r_k$. Hence $r_1=r_2=r_3=:r$ and $Z=(l_1+l_2+l_3)\wedge r$, which is impossible.
\end{enumerate}

\emph{Two full}, say $P_i,P_j$ full and $P_m$ of type $v=(v_1,v_2)$. Comparing the two components of the cross-part relation gives
\[
  r_i+r_j+v_1r_m=0,\qquad l_i+l_j+v_2r_m=0 .
\]
Substituting $l_j=l_i+v_2r_m$ and $r_j=r_i+v_1r_m$, and using $r_m\wedge r_m=0$,
\[
  Z=l_i\wedge r_i+l_j\wedge r_j+l_m\wedge r_m=(v_1l_i+v_2r_i+l_m)\wedge r_m ,
\]
which is impossible.

\medskip\noindent\emph{Circuits.} In the normal form of [L4, Lemma~1] with $\lambda=4$, the quadratic part of a linear gate $L_iR_i$ is $\ell(L_i)\wedge\ell(R_i)$, where $\ell(\cdot)$ is the linear part; this has rank at most $2$. By [L4, Lemma~2] and the paragraph after it, the last gate solves \eqref{eq:star} for $Q=\Span\{\pi_2(g_0),\dots,\pi_2(g_3)\}$, with $\pi_2(g_0)=P_0$ and $Z_1\in Q$. The first part of the theorem excludes this.
\end{proof}

\begin{remark}
Theorem~\ref{thm:main} needs neither $\dim Q=4$ nor $\operatorname{rank}P_i=2$. The reduction to that case in [L4, \S1, ``The case $\lambda=4$''] is therefore not needed, nor are [L4, Lemma~5], [L4, Proposition~7], the dimension statement of [L4, Lemma~4], or the identity $\pi_4(fh)=f_2\wedge h_2$. The split by $d_U$ in [L4] is replaced by the split by $|S|$. Lemma~\ref{lem:8p} absorbs both the case ``$d_U=3$, $R$ full'' (through Case~A) and every case that [L4] closes with Lemma~8.
\end{remark}

\section{Correspondence with the Lean development}\label{sec:table}

All names are in the namespace \lean{UnrestrictedBooleanMul.LowHalfMul3} of the repository \texttt{lowhalf-mul3}, at commit \texttt{4626be5} (Lean 4.35.0-rc3). The Lean statements use $\Phi(\omega)=(\iota_{e_{a_2}}\omega,\iota_{e_{b_2}}\omega)$ (\lean{Phi}), 2-forms as symmetric zero-diagonal $6\times6$ matrices over $\F$ (\lean{TwoForm6}), and $l\wedge r$ as \lean{wedge l r}. Each finite fact is checked by plain \lean{decide}, which the kernel replays; the table records which.

\begin{center}
\small
\begin{tabularx}{\textwidth}{@{}X l l@{}}
\toprule
Paper step & Lean theorem & File \\
\midrule
Statement of Theorem~\ref{thm:main} & \lean{StarFree} (definition), \lean{starFree} & \lean{Assembly}, \lean{MainTheorem} \\
Circuits: $\lambda=4$ excluded & \lean{lambdaCase_four_of_starFree}, \lean{lambdaCase_four} & \lean{Assembly}, \lean{MainTheorem} \\
\quad Lemma~2 of [L4] ($V_4=V_4(Q)$) & \lean{wireSpace_eq_V4} & \lean{LemmaTwo} \\
\quad Lemma~1 of [L4] (normal form) & \lean{normal_form_five} & \lean{NormalForm} \\
$\pi_2(LR)=\ell(L)\wedge\ell(R)$ for affine $L,R$ & \lean{pi2_mul_affine} & \lean{Product} \\
$\pi_2(z_1)=Z_1$, $\pi_2(a_0b_0)=P_0$ & \lean{pi2_lowMul3_one}, \lean{pi2_a0b0} & \lean{Product} \\
Bridge \eqref{eq:bridge} & \lean{deriv_add_deriv_zero_eq} & \lean{Contraction} \\
Lemma~\ref{lem:8p} (Lemma~8$'$) & \lean{starFree_of_goodDir} & \lean{LineDerivative} \\
\quad $D_ep=c_p+\gamma_p\ell$ on $V_4(Q)$ & \lean{deriv_line_of_mem} & \lean{LineDerivative} \\
\quad normalisation $\gamma_f=0$ & inside \lean{starFree_of_goodDir} & \lean{LineDerivative} \\
\quad product rule; $D_ez_2=\delta_e$ & \lean{deriv_mul}; \lean{PaperPair.deriv_z2} & \lean{LineDerivative} \\
\quad cases $c=1$ and $c=0$ & \lean{line_core} & \lean{LineDerivative} \\
\quad $\delta_e\notin\Span\{1,\ell\}\Rightarrow H\equiv1$ (\lean{decide}) & \lean{line_finite_b0}, \lean{_a0}, \lean{_a0b0} & \lean{LineDerivative} \\
\quad degree-4 part (the carry) & \lean{not_degreeLE_three_of_star} & \lean{LineDerivative} \\
\quad hypothesis ``$\iota_e(Q)\subseteq\langle\ell\rangle$'' & \lean{GoodDir}; \lean{goodDir_of_gens} & \lean{LineDerivative}, \lean{MainTheorem} \\
Lemma~8 of [L4] (case $\ell=0$) & \lean{derivative_lemma_paper}; \lean{goodDir_kill} & \lean{Contraction}, \lean{MainTheorem} \\
Case~A (Lemma~\ref{lem:A}) & \lean{goodDir_caseA} & \lean{MainTheorem} \\
\quad Lemma~6 of [L4] (\lean{decide}, $64\times64$) & \lean{lemma6} & \lean{CrossPart} \\
$\Phi(P_0)=0$, $\iota_e(P_0)=0$ & \lean{phi_P0}, \lean{iota_P0} & \lean{MainTheorem} \\
$\Phi(Z)=0$, $Z\in\Lam(U)$ & \lean{phi_Z}, \lean{Z_rows} & \lean{MainTheorem} \\
$Z$ is not of the form $l\wedge r$ (Pfaffian) & \lean{wedge_ne_Z} & \lean{LemmaThree} \\
Decomposition $Z_1=\sum c_kP_k$ & \lean{goodDir_main} & \lean{MainTheorem} \\
Split by $|S|$ & \lean{decomposition_cases} & \lean{MainTheorem} \\
\quad $|S|\le1$ & \lean{none_false}, \lean{single_false} & \lean{MainTheorem} \\
\quad $|S|=2$ & \lean{pair_zero} & \lean{MainTheorem} \\
\quad Lemma~3 of [L4], the two instances used (\lean{decide}) & \lean{lemma3}, \lean{lemma3_finite_Z1}, \lean{lemma3_finite_Z1P0} & \lean{LemmaThree} \\
\quad $|S|=3$ & \lean{triple_core} & \lean{MainTheorem} \\
\quad\quad a cross part is $0$ $\Rightarrow$ Case~A & inside \lean{triple_core} & \lean{MainTheorem} \\
Lemma~11 of [L4] (types), normal forms & \lean{types}, \lean{types_finite}, \lean{wedge_lincomb} & \lean{Types} \\
\quad cross parts in normal form & \lean{rows_v}, \lean{rows_full}, \lean{row_sum} & \lean{Types}, \lean{MainTheorem} \\
\quad contractions $\iota_e(l\wedge r)=l(e)\,r$ for type $v$ & \lean{iota_dirA2_v}, \lean{iota_dirB2_v}, \lean{iota_dirA2B2_v} & \lean{Types} \\
\quad $\gamma=0$ for type $v$, $1$ for full & \lean{gamma_v}, \lean{gamma_full} & \lean{Types} \\
Parity (one or three full) & inside \lean{triple_core} & \lean{MainTheorem} \\
None full, (a)--(c) (\lean{decide} on 6 bits) & \lean{zeroFull}, \lean{zeroFull_finite} & \lean{MainTheorem} \\
Two full & \lean{twoFull} & \lean{MainTheorem} \\
\bottomrule
\end{tabularx}
\end{center}

\paragraph{Two differences from the text.} Lemma~3 of [L4] is used only for $Z\in\{Z_1,Z_1+P_0\}$, and in Lean it is checked in Pfaffian form by \lean{decide}: if all fifteen $4\times4$ Pfaffians of $Z+l\wedge r$ vanish, then $\Phi(l\wedge r)=0$, for all $64\times64$ pairs $(l,r)$. Lemma~6 of [L4] is likewise a \lean{decide} over all pairs $(l,r)$, for all forms $l\wedge r$ including $0$. The human arguments for both lemmas in [L4] are not formalized.

\paragraph{Status.} Every step above is proved in Lean. The whole development builds with \lean{lake build} (1741 jobs) and contains no \lean{sorry}, \lean{native_decide}, \lean{ofReduceBool} or \lean{decide +kernel}. The audit file \lean{AxiomAudit.lean} reports the axioms \lean{propext}, \lean{Classical.choice}, \lean{Quot.sound} for all 130 theorems it lists. The text above was written by an AI agent from the Lean proof and has not been reviewed by a human.

\begin{thebibliography}{9}
\bibitem{L4} Bishara, \emph{A human proof of the four-linear-gate case for the $3\times3$ low-half multiplier}, \texttt{lambda4\_proof.tex} (11 October 2026, revised).
\bibitem{morse} G.~Morse, \emph{Unrestricted Boolean multiplicative complexity of four-term binary polynomial multiplication}, arXiv:2608.30238.
\end{thebibliography}

\end{document}
